% Part: first-order-logic
% Chapter: syntax-and-semantics
% Section: negative-free-logic

\documentclass[../../../include/open-logic-section]{subfiles}

\begin{document}

\olfileid{frl}{syn}{pfl}

\olsection{Positive Free Logic}

The guiding idea of positive free logic is that atomic claims do not have
existence committments. In natural language the following may seem true:

\begin{enumerate}
  \item Pegasus is Pegasus.
  \item Pegasus is a winged horse.
  \item Some children are waiting for Santa Claus.
  \item The round square is a square.
\end{enumerate}

If they are true, however, they should not entail that Pegasus, 
Santa Claus or the round square exist. Moreover, the following seem 
false:

\begin{enumerate}
  \item Pegasus is Santa Claus.
  \item Santa Claus is a winged horse.
  \item Some children are waiting for Pegasus.
\end{enumerate}

Accordingly, positive free logic allows that, where $a$ is empty in a 
!!{structure} $\Struct M$ and $F$ some predicate, $\Atom{F}{a}$ may be 
either true or false in $\Struct M$. So $\Atom{F}{a}$ does not 
entail $\lfrexists a$. Moreover, it allows that predicates behave 
differently with different empty !!{constant}s: for instance, where 
$a$ and $b$ are empty in a !!{structure} $\Struct M$ and $F$ some 
predicate, we may have $\Atom{F}{a}$ true and $\Atom{F}{b}$ false
in $\Struct M$. We will have $\eq[a][a]$ and $\eq[b][b]$ but 
$\eq[a][b]$ may be either true or false in a given $\Struct M$. In short,
the truth of falsity of atomic !!{formula}s with empty !!{constant}s is
\emph{non-trivial} in positive free logic.

\begin{explain}
To ensure atomic !!{formula}s with empty !!{constant}s have non-trivial
truth conditions, !!{structure}s of positive free logic use \emph{non-existing
objects}. That is, it gives empty !!{constant}s referents. These referents
can be part of the extension of predicates and are used to determine
the truth of falsity of claim using these !!{constant}s. However, these
referents are not considered as existing: they are not in the domain
over which $\lforall$ and $\lexists$ operate. 
\end{explain}

\end{document}
